% Section 2 -- The SPHERE-SPACE bridge
% Requires: \usepackage{graphicx}, \usepackage{booktabs}
% Placeholder citation keys: SPACE2025, SPHERE2025 (replace with the real ones)

\section{Bridging Pricing Design and Runtime Operation}
\label{sec:bridge}

The bridge consists of two complementary extensions. On the one hand, SPHERE now assigns each pricing a unique and stable identifier, and exposes its published versions through a public API. On the other hand, SPACE uses this identifier to link to a pricing in SPHERE and detect the release of new versions. Whenever a new version is detected, SPACE migrates the affected contracts and applies the changes to the host application. Once both tools are linked --a one-time configuration--, no further manual intervention is required.

Fig.~\ref{fig:bridge} depicts the resulting workflow using TomatoMeter, a fictitious productivity app whose features and usage limits are governed by a pricing managed through SPACE. The workflow comprises seven steps, from the moment a stakeholder edits this pricing in SPHERE until TomatoMeter reflects the change. The first two are performed by the stakeholder in SPHERE, while the remaining five are carried out automatically by SPACE and TomatoMeter.

\begin{figure}[t]
  \centering
  \includegraphics[width=\textwidth]{figures/bridge-workflow.png}
  \caption{Pricing update workflow enabled by the extensions of SPHERE and SPACE. The dashed arrow represents the periodic polling performed by SPACE; solid arrows represent the remaining requests and data flows. Numbers identify the steps described in Sec.~\ref{sec:bridge}.}
  \label{fig:bridge}
\end{figure}

\subsubsection{Pricing design and release (1--2).}
SPHERE~\cite{SPHERE2025} provides a pricing editor (1) where stakeholders can either edit a pricing visually, without dealing with its YAML serialization, or modify the iPricing directly. Once the changes are ready, the stakeholder publishes them as a new version of the pricing (2), which is stored in its version history alongside previous ones. To make versions consumable by external services, SPHERE assigns each pricing a stable identity, exposed through a permanent link (\texttt{/p/\{pricingId\}}) that survives renamings and relocations of the pricing. Behind this link, a read-only public API offers two operations: (i) a \emph{manifest} that lists the published versions of the pricing, each with an immutable identifier, its version label, its release date and a content hash; and (ii) the download of the iPricing of any of these versions. The content hash is a SHA-256 fingerprint computed over a canonical form of the iPricing, which allows consumers to detect any functional change in a version regardless of formatting.

\subsubsection{Linking SPACE to SPHERE.}
Instead of uploading an iPricing, the provider creates a SPACE service linked to the permanent link of the pricing. In this one-time configuration, the provider also selects how frequently SPACE must check SPHERE for updates (every 1, 5 or 15 minutes) and the \emph{synchronization policy} that determines how new versions affect existing subscriptions (see Table~\ref{tab:policies}). From then on, the pricing of the service can only evolve through SPHERE: manual uploads are rejected to prevent both tools from diverging.

\subsubsection{Version detection and import (3--4).}
The linked service activates a synchronization worker in SPACE, which periodically retrieves the manifest of the pricing through its permanent link (3) and selects a \emph{target version} according to the configured policy: either the latest published version, ordered by release date, or a specific version chosen by the provider. We opted for a pull-based design so that SPHERE does not need to know, authenticate or reach its consumers; it only exposes public data. Moreover, if SPHERE becomes unavailable, SPACE keeps operating with its local copies and retries in the next cycle. Only when a new target is detected does SPACE download its iPricing (4), check that its content matches the hash advertised in the manifest, and validate it before storing it as a new pricing version of the service. Imported versions are treated as immutable: if the content behind an already imported version changes, synchronization is blocked until a new version is published. Similarly, policies that follow the latest version never move a service back to an older one, even if the current target is unpublished.

\subsubsection{Contract migration (5).}
Once the target is imported, SPACE reconciles the subscriptions --i.e. contracts-- bound to the service according to the configured policy (Table~\ref{tab:policies}). Policies prefixed with \texttt{all} migrate every existing contract to the target version, whereas those prefixed with \texttt{new} preserve the version of existing contracts. In both cases, any new subscription is forced to use the target version, regardless of the version requested by the client.

\begin{table}[t]
  \centering
  \caption{Synchronization policies supported by SPACE.}
  \label{tab:policies}
  \begin{tabular}{llll}
    \toprule
    \textbf{Policy} & \textbf{Existing contracts} & \textbf{New contracts} & \textbf{Target version} \\
    \midrule
    \texttt{all\_last} & Migrated & Target version & Latest published version \\
    \texttt{new\_last} & Version kept & Target version & Latest published version \\
    \texttt{all\_pick} & Migrated & Target version & Version selected by provider \\
    \texttt{new\_pick} & Version kept & Target version & Version selected by provider \\
    \bottomrule
  \end{tabular}
\end{table}

When a contract is migrated, SPACE first tries to preserve its plan, add-ons and quantities. However, the new version may render this configuration invalid, e.g., if the plan or an add-on was removed. In such cases, the contract falls back to the cheapest plan of the new version that yields a valid subscription without add-ons. To avoid leaving the system in an inconsistent state, all contracts are prevalidated before any change is persisted: if a single contract cannot be migrated, none is. Migration also preserves billing dates and the subscription state. Usage levels of limits that remain in the new version keep their values, even if they exceed a reduced limit, whereas new limits start from their default values. If the semantics of a limit change in an incompatible way (e.g., its period or unit), migration is blocked and the service is flagged until the pricing is corrected. Every migration is recorded in the history of the contract, together with the previous plan and add-ons and the reason for the change, so that it can be audited. Finally, to cope with concurrent usage updates and partial failures, migrations rely on conditional writes and resume from where they stopped if interrupted.

\subsubsection{Notification and adaptation (6--7).}
After applying the changes, SPACE invalidates the cached pricings, contracts and feature evaluations affected by the synchronization, and broadcasts a \emph{pricing activated} event through WebSockets (6). Both the frontend and the backend of TomatoMeter are subscribed to these events through the SPACE SDKs, so the host application adapts without being redeployed (7): the frontend updates its components, while the backend updates the business logic that depends on the pricing. From then on, every feature evaluation is performed against the new version of the pricing.

Throughout this process, SPACE exposes the synchronization status of each service (e.g., up to date, applying changes, blocked or degraded), the version currently applied and the number of migrated contracts. This allows stakeholders to confirm that their changes have reached the running system, closing the loop without involving a developer.
